\chapter{Reproducibility and Artifact Evaluation}
\label{ch:artifact}

\section{Reproducing the Build}
\label{sec:artifact-build}

MATHLIB5 ships as a single cloneable repository with a deterministic
startup sequence (cf.\ AGENTS.md). To reproduce the verified artifacts:

\begin{lstlisting}[language=BashX,caption={Reproducing the MATHLIB5
  verification chain.},label=lst:repro]
git clone https://github.com/SNAPKITTYWEST/mathlib5 && cd mathlib5
npm ci
npm run verify:all        # plasma gate + P/NP proofs + skill seals
npm run context:bootstrap # load latest memories into local index
./scripts/receipt.sh      # emit the WORM-anchored verification receipt
\end{lstlisting}

A successful run produces a receipt whose SHA-256 is the artifact under
discussion in \cref{ch:anchors}.

\section{The Verification Gate}
\label{sec:artifact-gate}

The gate is \emph{total and fail-loud}: if any of the ten boundaries
returns a non-certified result, the build aborts before any downstream
artifact is emitted. This guarantees that a published artifact is
accompanied by a complete certificate chain or does not exist at all.

\section{The Anchors Published With This Paper}
\label{sec:artifact-anchors}

This monograph is itself anchored. The following artifacts were generated
and are verifiable:

\begin{itemize}
  \item \texttt{fingerprint.txt} --- SHA-256 of the rendered PDF.
  \item \texttt{pubkey.txt} --- the Ed25519 public key (Plasma Gate).
  \item \texttt{signature.txt} --- Ed25519 signature over the fingerprint.
  \item \texttt{opreturn.txt} --- the \texttt{MATH5:} Bitcoin
        \texttt{OP\_RETURN} payload (see \cref{ch:anchors}).
  \item \texttt{anchor\_meta.json} --- machine-readable anchor metadata.
\end{itemize}

Verification is performed by \texttt{verify\_anchors.py}, which checks
(1) that the recomputed SHA-256 equals \texttt{fingerprint.txt}, (2) that
the Ed25519 signature validates under \texttt{pubkey.txt}, and (3) that the
\texttt{OP\_RETURN} payload's hash matches the fingerprint. All three must
pass for the paper to be considered authentic.

\section{Evaluator Checklist}
\label{sec:artifact-checklist}

For artifact evaluation we recommend the following acceptance criteria:

\begin{enumerate}
  \item The repository clones cleanly and \texttt{npm ci} succeeds on a
        standard Linux environment.
  \item \texttt{npm run verify:all} exits zero.
  \item \texttt{./scripts/receipt.sh} emits a receipt whose hash matches
        the published fingerprint.
  \item \texttt{python verify\_anchors.py} reports \texttt{ALL CHECKS
        PASSED}.
\end{enumerate}

These criteria make the contribution \emph{externally reproducible} without
trusting the authors' hardware.
